%!TEX root = ../template.tex
%%%%%%%%%%%%%%%%%%%%%%%%%%%%%%%%%%%%%%%%%%%%%%%%%%%%%%%%%%%%%%%%%%%%
%% F1_introduction.tex
%% NOVA thesis document file
%%
%% Chapter with lots of dummy text
%%%%%%%%%%%%%%%%%%%%%%%%%%%%%%%%%%%%%%%%%%%%%%%%%%%%%%%%%%%%%%%%%%%%

\typeout{NT FILE F1_introduction.tex}%


%%%%%%%%%%%%%%%%%%%%%%%%%%%%%%%%%%%%%%%%%%%%%%%%%%
\chapter{Introdução}
\label{cha:introducao}
%%%%%%%%%%%%%%%%%%%%%%%%%%%%%%%%%%%%%%%%%%%%%%%%%%

%%%%%%%%%%%%%%%%%%%%%%%%%%%%%%%%%%%%%%%%%%%%%%%%%%
\section{Enquadramento e Motivação}
\label{sec:enquadramento_motivacao}
%%%%%%%%%%%%%%%%%%%%%%%%%%%%%%%%%%%%%%%%%%%%%%%%%%

Os \gls{ITS} constituem um dos mais importantes passos tecnológicos das cidades inteligentes, ao integrarem sistemas de comunicação, sensores, algoritmos inteligentes e infraestrutura digital de apoio à mobilidade. Estes sistemas emergem como resposta ao crescimento exponencial do tráfego urbano, à necessidade de reduzir acidentes e emissões, e à procura de soluções sustentáveis e eficientes de transporte \cite{zeadally_vehicular_2020}. A convergência entre tecnologias de comunicação veicular e computação distribuída tem impulsionado o desenvolvimento de aplicações avançadas, tais como gestão inteligente de semáforos, deteção cooperativa de perigos, coordenação de trajetórias em interseções, coordenação de frotas e monitorização ambiental, suportadas por veículos conectados e infraestrutura cooperativa.  

Um dos elementos mais críticos destes ecossistemas é a capacidade de comunicação rápida, fiável e escalável entre veículos. A eficácia de um sistema \gls{ITS} depende fortemente de como a informação é trocada, processada e atualizada entre os diversos participantes: veículos, \gls{RSUs} e centros de controlo. Contudo, as \gls{VANETs} caracterizam-se por terem desafios singulares: alta mobilidade, conectividade intermitente, densidade variável e topologias altamente dinâmicas. Tais características tornam ineficiente o paradigma de comunicação tradicional baseado em endereços IP, concebido originalmente para redes fixas e estáveis.  

Nas arquiteturas IP, a comunicação é orientada ao nó de destino: os nós da rede encaminham pacotes com base em endereços e rotas, independentemente do conteúdo transportado. Em ambientes veiculares, essa dependência de endereços fixos e de rotas estabelecidas resulta em latências elevadas, falhas frequentes de encaminhamento e baixa eficiência na disseminação de dados, especialmente em cenários de alta mobilidade ou desconexões momentâneas. Além disso, a necessidade constante de manutenção de tabelas de encaminhamento e reconfiguração de caminhos compromete a escalabilidade da rede.

Importa distinguir claramente dois conceitos frequentemente relacionados, mas distintos: a disseminação de informação e a sincronização de informação. A disseminação diz respeito ao processo de propagação de dados através da rede, garantindo que a informação chega a múltiplos nós, frequentemente de forma oportunista ou baseada em critérios espaciais ou temporais. Já a sincronização de informação visa assegurar que diferentes nós mantêm uma visão coerente e atualizada de um conjunto de dados partilhado, identificando e reconciliando versões divergentes ao longo do tempo.

Neste contexto, o paradigma \gls{NDN} surge como uma alternativa disruptiva à arquitetura IP. O \gls{NDN} inverte o modelo clássico: em vez de focar em “onde” estão os dados, concentra-se naquilo que se pretende obter. Assim, a comunicação é centrada nos conteúdos. Um veículo solicita informações por nome (por exemplo, “/cidade/avenida/prioridade/alerta”) sem conhecer o endereço IP do nó que contém informações. O resultado é uma rede mais flexível, segura e eficiente, que aproveita o \textit{in-network caching}, a replicação estratégica de conteúdos em múltiplos nós da rede (redundância natural dos dados, resultante da replicação oportunista de conteúdos em múltiplos nós através de mecanismos de caching distribuído) e o encaminhamento por nome para melhorar a entrega da informação em ambientes dinâmicos.

O \gls{NDN} tem demonstrado benefícios significativos em \gls{VANETs}, aumentando a robustez da comunicação, reduzindo o tráfego de controlo e permitindo que dados críticos, como avisos de acidente, alertas meteorológicos ou relativos a congestionamentos sejam rapidamente disseminados, mesmo com ligações intermitentes \cite{da_silva_realization_2022}. Além disso, a filosofia centrada em conteúdo alinha-se com o conceito de \textit{informação de interesse público}, essencial nos \gls{ITS}, onde o objetivo é garantir que a informação relevante chegue a quem dela necessita, no momento certo, independentemente de quem a produziu.  

Contudo, a disseminação e sincronização eficiente de dados em redes V-NDN continuam a ser um desafio complexo. A sincronização, neste contexto, refere-se ao processo pelo qual múltiplos nós da rede mantêm uma visão coerente e atualizada de um conjunto partilhado de dados distribuídos, garantindo que todos os participantes conhecem as versões mais recentes da informação produzida por cada membro do grupo. Este mecanismo é fundamental em aplicações veiculares cooperativas, onde veículos precisam de partilhar continuamente dados como posição, velocidade, alertas de segurança ou estado do tráfego, e onde qualquer nó deve conseguir identificar e obter rapidamente informação em falta sem sobrecarregar a rede com pedidos redundantes.
Em cenários urbanos ou rodoviários de elevada mobilidade, os veículos entram e saem constantemente da área de sincronização, gerando mudanças rápidas na topologia e dificultando a manutenção de estados partilhados. Protocolos de sincronização convencionais, como o SVSync, foram concebidos para redes fixas ou de baixa mobilidade, e não consideram as limitações de conectividade intermitente, latência variável ou restrições geográficas típicas das redes veiculares. A elevada mobilidade provoca três problemas específicos: primeiro, a janela temporal de contacto entre veículos é frequentemente inferior a 10 segundos, insuficiente para completar múltiplas rondas de sincronização; segundo, a composição do grupo de sincronização muda constantemente, dificultando a convergência para um estado global consistente; terceiro, o overhead de manutenção das estruturas de sincronização (como vetores de estado) pode crescer rapidamente em ambientes densos, comprometendo a escalabilidade do sistema.
A investigação na área destaca diversas lacunas ainda não resolvidas:

\begin{itemize}
    \item \textbf{Inadequação à alta mobilidade:} O \gls{SVSync} tradicional assume estabilidade temporal na rede, o que não se verifica em \gls{VANETs}, onde os nós mudam de posição e de vizinhança constantemente \cite{noauthor_ndn-waze_nodate}.
    \item \textbf{Ausência de estrutura hierárquica:} Faltam mecanismos de clusterização dinâmica que permitam organizar os veículos em grupos cooperativos, otimizando a sincronização local antes da disseminação global.
    \item \textbf{Escalabilidade limitada:} As abordagens existentes não exploram soluções que minimizem redundâncias de mensagens nem levam em consideração os diferentes níveis de densidade veicular.
    \item \textbf{Integração experimental insuficiente:} Observa-se uma escassez de implementações realistas do \gls{SVSync} adaptado a \gls{VANETs} no simulador \textit{ndnSIM}, dificultando a avaliação quantitativa e a reprodutibilidade dos resultados.
\end{itemize}

Torna-se, portanto, evidente a necessidade de um modelo de sincronização adaptativo e hierárquico, capaz de lidar com a volatilidade das conexões e de explorar a estrutura geográfica e temporal dos ambientes veiculares. Nestes cenários, a elevada taxa de entrada e saída de nós da rede provoca inconsistências frequentes nos estados sincronizados, aumentando o overhead de controlo e comprometendo a convergência temporal do protocolo.

A presente dissertação propõe uma camada de sincronização hierárquica, como resposta direta às limitações identificadas, inspirada nos princípios do SVSync, mas especificamente adaptada às particularidades das redes \gls{V-NDN}. Esta solução, utiliza clusterização hierárquica para otimizar a disseminação e adapta-se dinamicamente às variações da topologia. O resultado é um modelo mais eficiente e consistente de partilha de informação, que contribui para a fiabilidade dos sistemas cooperativos de transporte.  

Do ponto de vista prático, o impacto deste tipo de solução é vasto. Sistemas de coordenação de frotas, gestão de interseções, condução autónoma colaborativa e resposta rápida a emergências dependem da partilha e sincronização de informação em tempo real. Assim, qualquer avanço na eficiência dos mecanismos de sincronização traduz-se diretamente em melhorias práticas, por exemplo na disseminação atempada de alertas de travagem de emergência, ou na partilha coordenada de informação sobre o estado do tráfego para apoio à gestão inteligente de rotas, contribuindo para maior segurança rodoviária, redução de congestionamentos e melhor eficiência energética.

%%%%%%%%%%%%%%%%%%%%%%%%%%%%%%%%%%%%%%%%%%%%%%%%%%
\section{Objetivos}
\label{sec:objetivos}
%%%%%%%%%%%%%%%%%%%%%%%%%%%%%%%%%%%%%%%%%%%%%%%%%%

O objetivo central desta dissertação é conceber, implementar e avaliar uma solução de sincronização eficiente e escalável para redes veiculares baseadas em \gls{NDN}. Especificamente, pretende-se:

\begin{itemize}
    \item Analisar os principais mecanismos de sincronização do \gls{NDN}, com foco no \gls{SVSync}, identificando as suas limitações quando aplicados a ambientes de elevada mobilidade.
    \item Propor uma versão adaptada do \gls{SVSync}, concebida para \gls{VANETs}, integrando fatores como mobilidade e volumes de dados variáveis.
    \item Desenvolver o algoritmo de sincronização hierárquico, que organiza os veículos em múltiplas camadas (como \textit{lanes} e \textit{pivots}), reduzindo redundâncias e otimizando a partilha de informação.
    \item Implementar a solução no ambiente de simulação \textit{ndnSIM}, assegurando compatibilidade e capacidade de reprodução de cenários realistas.
    \item Avaliar o desempenho da proposta em termos de latência, taxa de sincronização e escalabilidade, comparando-a com a versão clássica do \gls{SVSync}.
\end{itemize}

%%%%%%%%%%%%%%%%%%%%%%%%%%%%%%%%%%%%%%%%%%%%%%%%%%
\section{Contribuições}
\label{sec:contribuicoes}
%%%%%%%%%%%%%%%%%%%%%%%%%%%%%%%%%%%%%%%%%%%%%%%%%%

Este trabalho oferece um conjunto de contributos científicos e tecnológicos que reforçam a integração do paradigma \gls{NDN} em redes veiculares inteligentes:

\begin{itemize}
    \item \textbf{Modelo Hierárquico Adaptativo:} Desenvolvimento de uma arquitetura de sincronização em múltiplos níveis, onde veículos “pivô” atuam como agregadores temporários de estado, reduzindo o tráfego redundante e otimizando a convergência de dados.
    \item \textbf{Camada SVSync-V:} Proposta de uma camada de sincronização específica para ambientes veiculares, incorporando métricas de localização na decisão de sincronização.
    \item \textbf{Integração Avançada com ndnSIM:} Implementação completa e funcional da solução no simulador \textit{ndnSIM}, com novos módulos e parâmetros configuráveis que permitem avaliar cenários de mobilidade complexa.
    \item \textbf{Avaliação Experimental Detalhada:} Comparação rigorosa entre a abordagem hierárquica proposta e o \gls{SVSync} convencional, quantificando os ganhos em eficiência e robustez.
\end{itemize}

%%%%%%%%%%%%%%%%%%%%%%%%%%%%%%%%%%%%%%%%%%%%%%%%%%
\section{Organização do Documento}
\label{sec:organizacao_documento}
%%%%%%%%%%%%%%%%%%%%%%%%%%%%%%%%%%%%%%%%%%%%%%%%%%

A dissertação encontra-se estruturada de forma a conduzir o leitor desde a contextualização teórica até à validação experimental da proposta:

\begin{itemize}
    \item O \textbf{Capítulo 1} introduz o enquadramento do tema, as motivações, os objetivos e as principais contribuições deste trabalho.  
    \item O \textbf{Capítulo 2} apresenta o estado da arte, abordando conceitos fundamentais de \gls{ITS}, redes veiculares, \gls{NDN} e protocolos de sincronização existentes.  
    \item O \textbf{Capítulo 3} descreve detalhadamente a arquitetura e o modelo de sincronização proposto, incluindo a organização hierárquica e o funcionamento dos algoritmos adaptativos.  
    \item O \textbf{Capítulo 4} discute a implementação prática no simulador \textit{ndnSIM}, explicando as estruturas de software, parâmetros configurados e metodologias de experimentação.  
    \item O \textbf{Capítulo 5} apresenta e analisa os resultados experimentais, comparando o desempenho da abordagem proposta com o \gls{SVSync} clássico e discutindo as melhorias alcançadas.  
    \item Por fim, o \textbf{Capítulo 6} sintetiza as conclusões do trabalho e propõe linhas de investigação futura, com vista à expansão da solução para cenários urbanos complexos e redes veiculares de larga escala.  
\end{itemize}


